\chapter{Tensor Products of Modules}\label{chap:tensor-products}

Direct sums place modules side by side. Tensor products solve a different
problem: they turn expressions that are additive in two variables into
ordinary homomorphisms from one abelian group. The price is that sidedness now
matters. We begin over an arbitrary ring and only later specialize to the
more familiar commutative case.

\section{Balanced Maps and Binary Tensor Products}\label{sec:binary-tensors}
\begin{definition}
Let \(M_R\) be a right \(R\)-module, \({}_RN\) a left \(R\)-module, and
\(A\) an abelian group. A map \(b:M\times N\to A\) is \emph{\(R\)-balanced}
if
\[
\begin{aligned}
b(m+m',n)&=b(m,n)+b(m',n),\\
b(m,n+n')&=b(m,n)+b(m,n'),\\
b(mr,n)&=b(m,rn).
\end{aligned}
\]
\end{definition}

The last equation is the balancing condition. It says that the scalar
inserted between the two variables may be moved across their interface. This
is precisely what makes a product such as \(m\otimes n\) meaningful when the
two module actions lie on opposite sides.

\begin{theorem}[Tensor-product universal property]\label{thm:binary-tensor}
There is an abelian group \(M\otimes_RN\), with an \(R\)-balanced map
\[
\tau:M\times N\longrightarrow M\otimes_RN,
\qquad (m,n)\longmapsto m\otimes n,
\]
such that every \(R\)-balanced map \(b:M\times N\to A\) has a unique
homomorphism \(\bar b:M\otimes_RN\to A\) satisfying
\(\bar b(m\otimes n)=b(m,n)\).
\end{theorem}
\[
\begin{tikzcd}[column sep=large, row sep=large]
M\times N \arrow[r, "\tau"] \arrow[dr, "b"'] &
M\otimes_R N \arrow[d, "\bar b"] \\
& A
\end{tikzcd}
\]
\begin{proof}
Let \(G\) be the free abelian group on symbols \([m,n]\), and quotient by
the subgroup generated by
\[
\begin{aligned}
[m+m',n]-[m,n]-[m',n],\qquad
[m,n+n']-[m,n]-[m,n'],\\
[mr,n]-[m,rn].
\end{aligned}
\]
Call the quotient \(M\otimes_RN\), and write \(m\otimes n\) for the class of
\([m,n]\). The displayed relations make \(\tau\) balanced. Given balanced
\(b\), the map from \(G\) sending \([m,n]\) to \(b(m,n)\) kills every
listed relation, so it descends to \(\bar b\). Since pure tensors generate
the quotient, this homomorphism is unique.
\end{proof}

If two groups satisfy this universal property, applying each universal
property to the other universal balanced map gives mutually inverse
homomorphisms. Thus \(M\otimes_RN\) is unique up to a \emph{canonical
isomorphism}. It is not a Cartesian product, and an element need not be a
single pure tensor; it is a finite sum of pure tensors.

\section{Bimodules and Tensor Actions}\label{sec:bimodules-tensors}
\begin{definition}
For rings \(R,S\), an \emph{\((R,S)\)-bimodule} \({}_RM_S\) is an abelian
group carrying a left \(R\)-action and a right \(S\)-action such that
\[
(rm)s=r(ms)
\qquad(r\in R,\ m\in M,\ s\in S).
\]
\end{definition}

Every ring is an \((R,R)\)-bimodule over itself. A left \(R\)-module is an
\((R,\mathbb Z)\)-bimodule and a right \(R\)-module is an
\((\mathbb Z,R)\)-bimodule. Bimodules retain an action on each outside of a
tensor product.

\begin{proposition}[Induced outside actions]\label{prop:tensor-actions}
Let \({}_SM_R\), \({}_RN_T\) be bimodules. Then \(M\otimes_RN\) is an
\((S,T)\)-bimodule under
\[
s(m\otimes n)=(sm)\otimes n,
\qquad
(m\otimes n)t=m\otimes(nt).
\]
\end{proposition}
\begin{proof}
For fixed \(s\), the map \((m,n)\mapsto(sm)\otimes n\) is balanced; indeed,
\[
s(mr)\otimes n=((sm)r)\otimes n=(sm)\otimes(rn).
\]
The universal property therefore gives a well-defined endomorphism. The
argument for \(t\) is analogous. The module axioms follow on pure tensors,
and the actions commute because
\[
(s(m\otimes n))t=(sm)\otimes(nt)=s((m\otimes n)t).
\]
\end{proof}

\section{Finite Associativity}\label{sec:tensor-associativity}
Associativity of tensor products is a canonical isomorphism, not literal
equality. Only finitely many tensor factors are considered.

\begin{theorem}[Associativity]\label{thm:tensor-associativity}
For bimodules \({}_SM_R\), \({}_RN_T\), and \({}_TP\), there is a canonical
isomorphism of left \(S\)-modules
\[
(M\otimes_RN)\otimes_TP
\ \cong\
M\otimes_R(N\otimes_TP),
\]
given by
\[
(m\otimes n)\otimes p\longmapsto m\otimes(n\otimes p).
\]
\end{theorem}
\begin{proof}
The proposed formula respects the \(R\)-balancing relation because
\[
(mr\otimes n)\otimes p
\longmapsto mr\otimes(n\otimes p)
=m\otimes(rn\otimes p).
\]
It respects the outer \(T\)-balancing relation because
\[
(m\otimes nt)\otimes p
\longmapsto m\otimes(nt\otimes p)
=m\otimes(n\otimes tp).
\]
Additivity is immediate, so the two universal properties produce a
well-defined homomorphism. The reversed formula
\[
m\otimes(n\otimes p)\longmapsto(m\otimes n)\otimes p
\]
satisfies the same checks. Both composites fix pure generators, hence are
the identity.
\end{proof}

\section{Tensor Products and Arbitrary Direct Sums}\label{sec:tensor-direct-sums}
This is the intended infinite-index construction: tensor product commutes
with arbitrary \emph{direct sums}, not arbitrary direct products. The
finite-support condition makes the following formula meaningful.

\begin{theorem}[Arbitrary direct sums]\label{thm:tensor-direct-sums}
For right \(R\)-modules \((M_i)_{i\in I}\) and a left \(R\)-module \(N\),
\[
\left(\bigoplus_{i\in I}M_i\right)\otimes_RN
\ \cong\
\bigoplus_{i\in I}(M_i\otimes_RN).
\]
For a right \(R\)-module \(M\) and left \(R\)-modules \((N_i)_{i\in I}\),
\[
M\otimes_R\left(\bigoplus_{i\in I}N_i\right)
\ \cong\
\bigoplus_{i\in I}(M\otimes_RN_i).
\]
\end{theorem}
\begin{proof}
For the first map, send \((m_i)\otimes n\) to the finite sum
\[
\sum_{i\in I}m_i\otimes n.
\]
It is balanced, so the tensor universal property makes it a homomorphism.
For each canonical injection
\[
\iota_i:M_i\longrightarrow\bigoplus_{i\in I}M_i,
\]
the balanced map \((m,n)\mapsto\iota_i(m)\otimes n\) induces
\[
M_i\otimes_RN\longrightarrow
\left(\bigoplus_{i\in I}M_i\right)\otimes_RN.
\]
The direct-sum universal property combines these maps into the reverse
homomorphism. The two maps are inverse on pure tensors and on the injected
summands, hence everywhere. The second isomorphism follows by the identical
argument with the variables reversed.
\end{proof}

\section{The Commutative-Ring Specialization}\label{sec:commutative-tensors}
Suppose now that \(R\) is commutative. Every left \(R\)-module becomes
canonically an \((R,R)\)-bimodule by defining \(mr=rm\). Thus
\(M\otimes_RN\) is again an \(R\)-module, and
\[
r(m\otimes n)=(rm)\otimes n=m\otimes(rn).
\]

\begin{proposition}[Symmetry]\label{prop:tensor-symmetry}
For \(R\)-modules \(M,N\), there is a canonical isomorphism
\[
M\otimes_RN\cong N\otimes_RM,
\qquad
m\otimes n\longmapsto n\otimes m.
\]
\end{proposition}
\begin{proof}
The map \((m,n)\mapsto n\otimes m\) is balanced because
\[
n\otimes(mr)=(rn)\otimes m.
\]
The universal property gives the displayed map; swapping the factors again
gives its inverse. Tensor products are symmetric only up to canonical
isomorphism, not by literal equality.
\end{proof}

\section{Finite Multilinear Tensor Products}\label{sec:multilinear-tensors}
For a finite list \(M_1,\ldots,M_n\) of \(R\)-modules and an \(R\)-module
\(A\), a map
\[
b:M_1\times\cdots\times M_n\longrightarrow A
\]
is \emph{\(R\)-multilinear} if it is additive and respects scalar
multiplication in each variable. Define
\[
M_1\otimes_R\cdots\otimes_RM_n
\]
to be the quotient of the free \(R\)-module on tuples
\((m_1,\ldots,m_n)\) by the additivity and scalar-linearity relations in each
slot. Write the class as \(m_1\otimes\cdots\otimes m_n\).

\begin{theorem}[Finite multilinear universal property]\label{thm:finite-tensors}
For every finite \(n\), \(R\)-module homomorphisms
\[
M_1\otimes_R\cdots\otimes_RM_n\longrightarrow A
\]
are in bijection with \(R\)-multilinear maps
\[
M_1\times\cdots\times M_n\longrightarrow A.
\]
For \(n=3\), the direct tensor product is canonically isomorphic to both
\[
(M_1\otimes_RM_2)\otimes_RM_3
\quad\text{and}\quad
M_1\otimes_R(M_2\otimes_RM_3).
\]
\end{theorem}
\begin{proof}
The quotient construction makes the universal tuple map multilinear. A
multilinear map from the product kills its defining relations and therefore
factors uniquely through the quotient, proving the bijection. For \(n=3\),
the map
\[
m_1\otimes m_2\otimes m_3
\longmapsto
(m_1\otimes m_2)\otimes m_3
\]
is well-defined by multilinearity and has the evident inverse on pure
tensors. Theorem~\ref{thm:tensor-associativity} gives the other
parenthesization. Induction on \(n\) proves the finite case.
\end{proof}

This justifies suppressing parentheses in finite tensor products. Repeated
associativity and symmetry also let every permutation of finitely many factors
induce a canonical isomorphism. No countable or arbitrary infinite tensor
product is being defined here.

\section{Tensoring Homomorphisms and Concrete Calculations}\label{sec:tensor-maps-calculations}
If \(f:M_R\to M'_R\) and \(g:{}_RN\to{}_RN'\) are module homomorphisms, then
\[
(f\otimes g)(m\otimes n)=f(m)\otimes g(n)
\]
defines a homomorphism \(f\otimes g:M\otimes_RN\to M'\otimes_RN'\).
Indeed, the displayed expression is balanced, so the universal property
proves well-definedness. On pure tensors,
\[
(f'\otimes g')\circ(f\otimes g)=(f'f)\otimes(g'g),
\qquad
1_M\otimes1_N=1_{M\otimes_RN}.
\]

Return to a commutative ring \(R\). The following calculations anchor the
abstract construction.

\begin{proposition}\label{prop:tensor-calculations}
There are canonical isomorphisms
\[
R\otimes_RM\cong M,
\qquad
R^m\otimes_RR^n\cong R^{mn}.
\]
If \(I\) is an ideal, then
\[
(R/I)\otimes_RM\cong M/IM.
\]
\end{proposition}
\begin{proof}
The balanced map \((r,m)\mapsto rm\) induces \(R\otimes_RM\to M\), with
inverse \(m\mapsto1\otimes m\). For finite free modules, \(e_i\otimes f_j\)
map to the \((i,j)\)-th standard vector; expanding both factors proves that
this map and its evident inverse are mutually inverse. Finally,
\((r+I,m)\mapsto rm+IM\) is balanced and well-defined because changing \(r\)
by an element of \(I\) changes the value by an element of \(IM\). Its inverse
is \(m+IM\mapsto(1+I)\otimes m\); balancing kills every generator \(im\).
\end{proof}

\begin{example}
Taking \(R=\mathbb Z\) and \(I=(m)\) gives
\[
\mathbb Z/m\mathbb Z\otimes_{\mathbb Z}\mathbb Z/n\mathbb Z
\cong
\mathbb Z/\gcd(m,n)\mathbb Z.
\]
Indeed, the preceding proposition identifies this with the quotient of
\(\mathbb Z/n\mathbb Z\) by the subgroup generated by \(m\).
\end{example}

\section{The Tensor--Hom Correspondence}
The universal property of a tensor product can be tested against arbitrary
abelian groups.  This produces the basic tensor--Hom correspondence without
assuming commutativity of the ring.

\begin{theorem}[Tensor--Hom]\label{thm:tensor-hom}
Let \(M_R\) be a right \(R\)-module and let \({}_RN\) be a left \(R\)-module.
Let \(A\) be an abelian group.  Give
\[
\operatorname{Hom}_{\mathbb Z}(M,A)
\]
the left \(R\)-module structure
\[
(r\cdot f)(m)=f(mr).
\]
Then there is an isomorphism of abelian groups
\[
\operatorname{Hom}_{\mathbb Z}(M\otimes_RN,A)
\cong
\operatorname{Hom}_R\!\left(N,\operatorname{Hom}_{\mathbb Z}(M,A)\right).
\]
\end{theorem}
\begin{proof}
Given \(u:M\otimes_RN\to A\), set
\[
\widehat u(n)(m)=u(m\otimes n).
\]
This is additive in \(n\).  Balancing gives
\[
\widehat u(rn)(m)=u(m\otimes rn)=u(mr\otimes n)
=(r\cdot\widehat u(n))(m),
\]
so \(\widehat u\) is \(R\)-linear.  Conversely, for an \(R\)-linear map
\(v:N\to\operatorname{Hom}_{\mathbb Z}(M,A)\), the rule
\[
(m,n)\longmapsto v(n)(m)
\]
is additive in each variable and balanced, since
\[
v(n)(mr)=(r\cdot v(n))(m)=v(rn)(m).
\]
It therefore induces \(M\otimes_RN\to A\).  Evaluating each composite on a
pure tensor shows that the two constructions are inverse.
\end{proof}

\begin{proposition}[Compatibility of Tensor--Hom]\label{prop:tensor-hom-compat}
The correspondence in Theorem~\ref{thm:tensor-hom} has the following
explicit compatibilities.  If \(f:M'\to M\) is a map of right modules and
\(u:M\otimes_RN\to A\), then \(u\circ(f\otimes1_N)\) corresponds to
\[
n\longmapsto\widehat u(n)\circ f.
\]
If \(h:A\to A'\) is a homomorphism of abelian groups, then \(h\circ u\)
corresponds to
\[
n\longmapsto h\circ\widehat u(n).
\]
Finally, if \(\ell:N'\to N\) is a map of left modules, then
\(u\circ(1_M\otimes\ell)\) corresponds to \(\widehat u\circ\ell\).
\end{proposition}
\begin{proof}
For example, the first map sends \(n\) and \(m'\) to
\[
u(f(m')\otimes n)=\widehat u(n)(f(m')).
\]
The second and third identities follow in exactly the same way by evaluating
on \(m\otimes n\) and \(m\otimes n'\), respectively.  The displayed
formulas also show that the resulting maps have the required module
linearity.
\end{proof}

\begin{proposition}[Commutative-ring Tensor--Hom]\label{prop:comm-tensor-hom}
If \(R\) is commutative and \(M,N,P\) are \(R\)-modules, then
\[
\operatorname{Hom}_R(M\otimes_RN,P)
\cong
\operatorname{Hom}_R\!\left(M,\operatorname{Hom}_R(N,P)\right)
\cong
\operatorname{Hom}_R\!\left(N,\operatorname{Hom}_R(M,P)\right).
\]
\end{proposition}
\begin{proof}
An \(R\)-linear map \(M\otimes_RN\to P\) is the same, by the tensor-product
universal property, as an \(R\)-balanced bilinear map \(M\times N\to P\).
Sending such a map \(b\) to
\[
m\longmapsto\bigl(n\longmapsto b(m,n)\bigr)
\]
gives the first isomorphism, with inverse obtained by evaluation.  Switching
the two variables gives the second.  Commutativity is what identifies the
two sidedness conventions and makes each Hom group an \(R\)-module.
\end{proof}

\begin{remark}
The preceding compatibility identities are often summarized by saying that
the Tensor--Hom correspondence is natural in its variables.  The explicit
sidedness in Theorem~\ref{thm:tensor-hom} remains essential when \(R\) is
not commutative.
\end{remark}

\section{Right Exactness}\label{sec:tensor-right-exactness}
Tensor product preserves cokernels, which is stronger than merely preserving
surjections.  Indeed, an exact sequence
\[
A\longrightarrow B\longrightarrow C\longrightarrow0
\]
says that \(C\cong B/\operatorname{im}(A)\), so tensoring preserves this
quotient/cokernel even though it may destroy injectivity of \(A\to B\).

\begin{theorem}[Right exactness]\label{thm:tensor-right-exact}
Let
\[
A_R\xrightarrow{u}B_R\xrightarrow{v}C_R\longrightarrow0
\]
be exact, and let \({}_RN\) be a left \(R\)-module. Then
\[
A\otimes_RN\xrightarrow{u\otimes1}
B\otimes_RN\xrightarrow{v\otimes1}
C\otimes_RN\longrightarrow0
\]
is exact.
\end{theorem}
\begin{proof}
We prove that \(v\otimes1\) is a cokernel of \(u\otimes1\).  Let \(X\) be an
abelian group and let
\[
F:B\otimes_RN\longrightarrow X
\]
satisfy \(F(u\otimes1)=0\).  Theorem~\ref{thm:tensor-hom} identifies \(F\)
with
\[
\widehat F:N\longrightarrow\operatorname{Hom}_{\mathbb Z}(B,X).
\]
By Proposition~\ref{prop:tensor-hom-compat}, the vanishing condition says
\[
\widehat F(n)\circ u=0
\]
for every \(n\in N\).  Since \(C\) is the cokernel of \(u\), each
\(\widehat F(n)\) factors uniquely as
\[
\widehat F(n)=G_n\circ v
\]
for a homomorphism \(G_n:C\to X\).  The assignment is additive because
\[
\bigl(G_{n+n'}-G_n-G_{n'}\bigr)\circ v=0,
\]
so uniqueness of factorization through the cokernel gives
\[
G_{n+n'}=G_n+G_{n'}.
\]
For the left \(R\)-action on \(\operatorname{Hom}_{\mathbb Z}(C,X)\),
\[
(r\cdot G_n)\circ v
=r\cdot(G_n\circ v)
=r\cdot\widehat F(n)
=\widehat F(rn).
\]
Uniqueness similarly gives \(G_{rn}=r\cdot G_n\).  Therefore
\[
n\longmapsto G_n
\]
is an \(R\)-linear map \(N\to\operatorname{Hom}_{\mathbb Z}(C,X)\).
Applying Theorem~\ref{thm:tensor-hom} backward gives a unique
\[
G:C\otimes_RN\longrightarrow X
\]
with \(G(v(b)\otimes n)=F(b\otimes n)\).  Hence
\[
F=G\circ(v\otimes1),
\]
and uniqueness follows because pure tensors generate.  Thus \(v\otimes1\)
is the cokernel of \(u\otimes1\), which is exactly the asserted exactness.
\end{proof}

\begin{corollary}[Tensoring a quotient]\label{cor:tensor-quotient}
For a map \(u:A_R\to B_R\) and a left \(R\)-module \({}_RN\), there is a
canonical isomorphism
\[
\bigl(B/\operatorname{im}u\bigr)\otimes_RN
\cong
\bigl(B\otimes_RN\bigr)/\operatorname{im}(u\otimes1).
\]
\end{corollary}
\begin{proof}
Apply Theorem~\ref{thm:tensor-right-exact} to the quotient map
\[
B\longrightarrow B/\operatorname{im}u.
\]
The two modules displayed are cokernels of the same map \(u\otimes1\), and
the unique map induced by the quotient is the required canonical
isomorphism.
\end{proof}

No injectivity statement is implied. Tensoring
\[
2\mathbb Z\hookrightarrow\mathbb Z
\]
with \(\mathbb Z/2\mathbb Z\) gives a zero map from a nonzero module. Later
homological algebra measures this failure; the next definition isolates when
the failure does not occur.

\section{Flat Modules}
\begin{definition}
A left \(R\)-module \(N\) is \emph{flat} if tensoring with \(N\) takes every
short exact sequence of right \(R\)-modules to a short exact sequence.
\end{definition}
Equivalently, tensoring with \(N\) preserves injections, since
Theorem~\ref{thm:tensor-right-exact} already gives right exactness.

\begin{proposition}\label{prop:projective-flat}
Every free module is flat, and every projective module is flat.  Thus
\[
\text{free}\Longrightarrow\text{projective}\Longrightarrow\text{flat}.
\]
\end{proposition}
\begin{proof}
If \(F\) is free on \(X\), then, for every right \(R\)-module \(A\),
\[
A\otimes_R F\cong\bigoplus_{x\in X}A.
\]
Consequently tensoring a short exact sequence of right modules with \(F\) is a direct sum of
short exact sequences and remains exact.  Thus free modules are flat.

By Theorem~\ref{thm:projective-characterizations}, a projective module \(P\)
is a direct summand of a free module \(F=P\oplus P'\).  Tensoring distributes
over this finite direct sum, so a tensor map involving \(F\) is the direct
sum of the corresponding tensor maps involving \(P\) and \(P'\).  If the
first is injective, each summand map is injective; if the middle sequence is
exact, exactness holds componentwise.  Since \(F\) is flat, \(P\) is flat.
\end{proof}
The converses need not hold in general.

\section{Exterior Powers and Exterior Algebra}\label{sec:exterior-algebra}
Finite tensor products represent multilinear maps.  Exterior powers impose
one further relation: a multilinear expression should vanish whenever two
inputs agree.  This packages alternating forms, oriented volume, and the
determinant into one construction.

Throughout this section \(R\) is a commutative ring and \(M\) is an
\(R\)-module.  A multilinear map
\[
f:M^k\longrightarrow N
\]
is \emph{alternating} if
\[
f(m_1,\ldots,m_k)=0
\]
whenever two arguments coincide.  This is the same alternating condition
used for determinants in \hyperref[sec:determinants]{Section~6.12}; it is
not to be confused with skew-symmetry in characteristic \(2\).

\begin{theorem}[Exterior-power universal property]\label{thm:exterior-power}
For every integer \(k\geq1\), there is an \(R\)-module
\[
\bigwedge_R^kM
\]
and an alternating multilinear map
\[
\omega_k:M^k\longrightarrow\bigwedge_R^kM,
\qquad
(m_1,\ldots,m_k)\longmapsto m_1\wedge\cdots\wedge m_k,
\]
such that every alternating multilinear map
\[
f:M^k\longrightarrow N
\]
has a unique linear factorization
\[
\widetilde f:\bigwedge_R^kM\longrightarrow N.
\]
\end{theorem}
\[
\begin{tikzcd}[column sep=large, row sep=large]
M^k \arrow[r, "\omega_k"] \arrow[dr, "f"'] &
\bigwedge_R^k M \arrow[d, "\widetilde f"] \\
& N
\end{tikzcd}
\]
\begin{proof}
By Theorem~\ref{thm:finite-tensors}, form the finite tensor power
\[
M^{\otimes k}
=M\otimes_R\cdots\otimes_RM.
\]
Let \(J\) be the submodule generated by all pure tensors having two equal
entries, and define
\[
\bigwedge_R^kM=M^{\otimes k}/J.
\]
The quotient of the universal multilinear tensor map is \(\omega_k\), and
it is alternating by construction.  If \(f\) is alternating, its associated
linear map
\[
M^{\otimes k}\longrightarrow N
\]
vanishes on every generator of \(J\), so it descends uniquely through the
quotient.  This gives the required factorization.  Any two modules with this
property are canonically isomorphic by applying their universal properties
to one another's universal alternating maps.
\end{proof}

\begin{proposition}[Alternating wedge identities]\label{prop:wedge-signs}
If two entries agree, then
\[
m_1\wedge\cdots\wedge m_k=0.
\]
Interchanging two entries changes sign; consequently
\[
m_{\sigma(1)}\wedge\cdots\wedge m_{\sigma(k)}
=\operatorname{sgn}(\sigma)\,
m_1\wedge\cdots\wedge m_k
\]
for every \(\sigma\in S_k\).
\end{proposition}
\begin{proof}
The first assertion is part of the construction.  Holding all other entries
fixed, alternation and multilinearity give
\[
0=\omega_k(\ldots,x+y,\ldots,x+y,\ldots).
\]
The repeated-\(x\) and repeated-\(y\) terms vanish, leaving the sum of the
two wedges with \(x\) and \(y\) interchanged.  Thus a transposition
contributes \(-1\).  The sign theory of Chapter~2 then gives the formula for
an arbitrary permutation.
\end{proof}

\noindent\textbf{Example.} If \(M=R^3\) has basis \(e_1,e_2,e_3\), then
\[
\bigwedge_R^2M
\]
has basis
\[
e_1\wedge e_2,\qquad e_1\wedge e_3,\qquad e_2\wedge e_3.
\]
Terms \(e_i\wedge e_i\) vanish, and reversing two basis vectors changes the
sign, so these are precisely the increasing-index wedges.

\begin{theorem}[Basis of an exterior power]\label{thm:exterior-basis}
If \(M\) is finite free with ordered basis
\[
(e_1,\ldots,e_n),
\]
then the wedges
\[
e_{i_1}\wedge\cdots\wedge e_{i_k},
\qquad
1\leq i_1<\cdots<i_k\leq n,
\]
form a basis of
\[
\bigwedge_R^kM.
\]
Hence
\[
\operatorname{rank}_R\bigl(\bigwedge_R^kM\bigr)=\binom nk.
\]
\end{theorem}
\begin{proof}
Multilinearity and Proposition~\ref{prop:wedge-signs} show that these wedges
span: expand every argument in the basis, discard repeated-index terms, and
reorder the rest.  To prove independence, fix an increasing \(k\)-tuple
\[
I=(i_1,\ldots,i_k).
\]
For \(m_j=\sum_ra_{rj}e_r\), define
\[
f_I(m_1,\ldots,m_k)=
\det(a_{i_rj})_{1\leq r,j\leq k}.
\]
The determinant theory of
\hyperref[sec:determinants]{Section~6.12} shows that \(f_I\) is alternating
and multilinear, so it induces a linear map from \(\bigwedge_R^kM\) to
\(R\).  On the displayed wedges it has value \(1\) for the tuple \(I\) and
\(0\) for every other increasing tuple.  Applying these maps to a linear
relation proves that every coefficient is zero.
\end{proof}

We set
\[
\bigwedge_R^0M=R,
\qquad
\bigwedge_R^1M=M.
\]
The basis theorem gives
\[
\bigwedge_R^kM=0
\qquad(k>n)
\]
when \(M\) is finite free of rank \(n\), and it shows that
\[
\bigwedge_R^nM
\]
is free of rank one.

\begin{proposition}[Induced exterior maps]\label{prop:exterior-maps}
An \(R\)-linear map
\[
T:M\longrightarrow N
\]
induces a unique linear map
\[
\bigwedge^kT:\bigwedge_R^kM\longrightarrow\bigwedge_R^kN
\]
satisfying
\[
\bigwedge^kT(m_1\wedge\cdots\wedge m_k)
=Tm_1\wedge\cdots\wedge Tm_k.
\]
Moreover,
\[
\bigwedge^k(\operatorname{id}_M)
=\operatorname{id}_{\bigwedge_R^kM},
\qquad
\bigwedge^k(S\circ T)
=\bigl(\bigwedge^kS\bigr)\circ\bigl(\bigwedge^kT\bigr).
\]
\end{proposition}
\begin{proof}
The map
\[
(m_1,\ldots,m_k)\longmapsto Tm_1\wedge\cdots\wedge Tm_k
\]
is alternating and multilinear.  By Theorem~\ref{thm:exterior-power}, it
therefore induces a unique map.  The two displayed identities agree
on every pure wedge, and pure wedges generate the relevant exterior powers.
\end{proof}

\begin{theorem}[Top exterior action]
Let \(V\) be \(n\)-dimensional over \(F\), and let \(T:V\to V\).  Then
\[
\bigwedge^nT
\]
acts on the one-dimensional space
\[
\bigwedge_F^nV
\]
by multiplication by \(\det(T)\).  Equivalently,
\[
Tv_1\wedge\cdots\wedge Tv_n
=\det(T)\,
\bigl(v_1\wedge\cdots\wedge v_n\bigr).
\]
\end{theorem}
\begin{proof}
Fix an ordered basis \((v_1,\ldots,v_n)\).  The linear functional on
\(\bigwedge_F^nV\) that takes \(v_1\wedge\cdots\wedge v_n\) to \(1\)
corresponds, by the exterior universal property, to a normalized alternating
\(n\)-linear form on \(V\).  By the uniqueness theorem for determinants in
\hyperref[sec:determinants]{Section~6.12}, its value on
\[
(Tv_1,\ldots,Tv_n)
\]
is \(\det(T)\).  Since the top exterior power is one-dimensional, this is
exactly the scalar by which \(\bigwedge^nT\) acts.
\end{proof}

This gives a later, equivalent characterization of the determinant already
developed from normalized alternating multilinear forms in
\hyperref[sec:determinants]{Section~6.12}.  Since
\[
\dim_F\bigwedge_F^nV=1,
\]
the determinant of \(T\) is the unique scalar \(\lambda\in F\) such that
\[
\bigwedge^nT
=\lambda\operatorname{id}_{\bigwedge_F^nV}.
\]
Equivalently, for every ordered basis \((v_1,\ldots,v_n)\) of \(V\),
\[
Tv_1\wedge\cdots\wedge Tv_n
=\det(T)\bigl(v_1\wedge\cdots\wedge v_n\bigr).
\]
Thus this top-exterior-power viewpoint describes the same invariant as the
Leibniz and alternating-multilinear descriptions, rather than defining a
new determinant.

For \(A\in M_n(F)\), regarded as an endomorphism of \(F^n\), the matrix of
\[
\bigwedge^nA
\]
with respect to the one-element basis
\[
e_1\wedge\cdots\wedge e_n
\]
of \(\bigwedge_F^nF^n\) is the \(1\times1\) matrix
\[
(\det A).
\]
Moreover, the composition law proved in Proposition~\ref{prop:exterior-maps}
makes determinant multiplicativity conceptually immediate:
\[
\bigwedge^n(S\circ T)
=\bigl(\bigwedge^nS\bigr)\circ\bigl(\bigwedge^nT\bigr)
\]
implies
\[
\det(S\circ T)=\det(S)\det(T).
\]
This is a retrospective explanation of the multiplicativity already proved
in Chapter~6, not a replacement for that earlier proof.

\section{The Exterior Algebra and Alternating Forms}\label{sec:exterior-forms}
The individual exterior powers fit together into a graded algebra.  This
uses an arbitrary direct sum of already constructed modules; it does
\emph{not} introduce an infinite tensor product.

\begin{definition}
The \emph{exterior algebra} of \(M\) is
\[
\bigwedge M
=\bigoplus_{k\geq0}\bigwedge_R^kM.
\]
For pure wedges, define the product by concatenation:
\[
\begin{aligned}
&(m_1\wedge\cdots\wedge m_p)
\wedge
(n_1\wedge\cdots\wedge n_q)\\
&\qquad =
m_1\wedge\cdots\wedge m_p
\wedge n_1\wedge\cdots\wedge n_q.
\end{aligned}
\]
\end{definition}

The right-hand side is alternating separately in the \(m\)'s and in the
\(n\)'s, so the two exterior universal properties make this a well-defined
bilinear map
\[
\bigwedge_R^pM\times\bigwedge_R^qM
\longrightarrow
\bigwedge_R^{p+q}M.
\]
Pure wedges generate each factor, so associativity follows from associativity
of concatenation on pure wedges.  The degree-\(k\) summand is the \(k\)-th
graded piece.

\begin{proposition}[Graded commutativity]
For homogeneous elements
\[
\alpha\in\bigwedge_R^pM,
\qquad
\beta\in\bigwedge_R^qM,
\]
one has
\[
\alpha\wedge\beta=(-1)^{pq}\beta\wedge\alpha.
\]
\end{proposition}
\begin{proof}
For pure wedges, moving the block of \(p\) entries past the block of \(q\)
entries requires \(pq\) transpositions.  Proposition~\ref{prop:wedge-signs}
gives the sign.  Bilinearity extends the identity to arbitrary homogeneous
elements.  The argument uses alternation itself, so it remains valid in
characteristic \(2\), where the displayed sign may equal \(1\).
\end{proof}

\begin{theorem}[Alternating forms as exterior powers of the dual]
Let \(V\) be a finite-dimensional vector space over \(F\).  There is a
natural isomorphism from
\[
\bigwedge_F^kV^\vee
\]
onto the space of alternating \(k\)-linear maps
\[
V^k\longrightarrow F.
\]
On a decomposable element it is given by
\[
\varphi_1\wedge\cdots\wedge\varphi_k
\longmapsto
\bigl((v_1,\ldots,v_k)\mapsto
\det(\varphi_i(v_j))_{1\leq i,j\leq k}\bigr).
\]
\end{theorem}
\begin{proof}
The displayed determinant is alternating and multilinear in the vectors, so
the exterior universal property makes the assignment well-defined and
linear.  If \((v_1,\ldots,v_n)\) is a basis of \(V\), then the wedges of its
dual basis indexed by increasing \(k\)-tuples map to the corresponding
coordinate alternating forms.  These forms are linearly independent and
span by evaluating an alternating form on basis tuples.  Thus the map is an
isomorphism.
\end{proof}

This identifies the algebra developed here with the pointwise algebra behind
differential forms.  On a smooth manifold, a differential \(k\)-form assigns
to each point \(p\) an alternating \(k\)-linear form on the tangent space
\[
T_pM,
\]
that is, an element of
\[
\bigwedge^k(T_pM)^\vee.
\]
The wedge product sends a \(p\)-form and a \(q\)-form at \(p\) to an element
of
\[
\bigwedge^{p+q}(T_pM)^\vee.
\]
No manifold theory is needed here; the point is that exterior algebra is the
linear-algebraic model that makes those geometric objects possible.


\section*{Exercises}
\begin{exercise}[Basic] Verify that multiplication \(R\times R\to R\) is
balanced when the first copy is the regular right module and the second the
regular left module.\end{exercise}
\begin{exercise}[Core] Give an \((R,R)\)-bimodule structure on \(R^n\), and
verify the induced actions of Proposition~\ref{prop:tensor-actions}.
\end{exercise}
\begin{exercise}[Core] Construct explicitly the associativity isomorphism
for \(\mathbb Z/2\mathbb Z\), \(\mathbb Z/3\mathbb Z\), and
\(\mathbb Z/4\mathbb Z\).\end{exercise}
\begin{exercise}[Core] Prove the second arbitrary-direct-sum isomorphism in
Theorem~\ref{thm:tensor-direct-sums} directly.\end{exercise}
\begin{exercise}[Core] Compute \((R/(a))\otimes_R(R/(b))\) for a commutative
ring \(R\).\end{exercise}
\begin{exercise}[Further] Verify the universal property of the three-factor
tensor product directly from the quotient construction.\end{exercise}
\begin{exercise}[Further, Challenging] Show that tensoring the inclusion
\(2\mathbb Z\hookrightarrow\mathbb Z\) with \(\mathbb Z/2\mathbb Z\) is not
injective.\end{exercise}
\begin{exercise}[Core] If \(M\) is free with basis \(e_1,e_2,e_3\), give a
basis for
\[
\bigwedge_R^2M.
\]
For a diagonal endomorphism of \(F^3\), compute its induced map on this
exterior power and recover its determinant from the top exterior power.
\end{exercise}
